\chapter{Описание модели}

В данной главе представлена стандартная модель реального делового цикла (Real Business Cycle, RBC) с дискретным временем. Экономика состоит из репрезентативного домохозяйства, репрезентативной фирмы и государства (которое для простоты отсутствует). Модель описывает поведение экономических агентов в условиях совершенной конкуренции, отсутствия внешних эффектов и полной предопределённости (детерминированная версия). Все переменные измеряются в реальном выражении.

\section{Экономическая среда}

Время дискретно, \(t = 0, 1, 2, \dots, \infty\). В каждом периоде домохозяйство принимает решения о потреблении \(C_t\) и предложении труда \(L_t\). Фирма производит однородный конечный продукт \(Y_t\) с использованием капитала \(K_t\) и труда \(L_t\) по технологии Кобба–Дугласа. Капитал накапливается в соответствии с законом движения, включающим амортизацию \(\delta \in (0,1)\). Технологический прогресс отсутствует, уровень совокупной факторной производительности (TFP) \(A\) постоянен.

\section{Задача домохозяйства}

Домохозяйство максимизирует дисконтированную сумму мгновенных полезностей, зависящих от потребления и труда:


\[
\max_{\{C_t, L_t\}_{t=0}^\infty} \sum_{t=0}^\infty \beta^t u(C_t, L_t),
\]

где \(\beta \in (0,1)\) – субъективный дисконт-фактор. Функция полезности задаётся в сепарабельном виде с постоянной эластичностью замещения (CRRA) по потреблению и линейным отрицательным вкладом труда:

\[
u(C_t, L_t) = \frac{C_t^{1-\sigma} - 1}{1-\sigma} - \varphi \frac{L_t^{1+\psi}}{1+\psi},
\]

где \(\sigma > 0\) – коэффициент относительного неприятия риска (обратная эластичность межвременного замещения), \(\varphi > 0\) – параметр, отражающий «ненависть к труду», \(\psi > 0\) – обратная эластичность предложения труда по Фришу.

Домохозяйство владеет начальным запасом капитала \(K_0\) и получает доход от сдачи капитала в аренду по ставке \(R_t\) и от продажи труда по ставке заработной платы \(w_t\). Доход используется на потребление и инвестиции \(I_t\). Бюджетное ограничение в каждом периоде имеет вид:

\[
C_t + I_t = w_t L_t + R_t K_t.
\]

Капитал изменяется согласно закону движения:

\[
K_{t+1} = I_t + (1-\delta)K_t.
\]

Домохозяйство выбирает \(\{C_t, L_t, K_{t+1}\}_{t=0}^\infty\), максимизируя полезность при заданных начальных условиях и бюджетных ограничениях.

\section{Задача фирмы}

Фирма производит продукт по технологии Кобба–Дугласа с постоянной отдачей от масштаба:

\[
Y_t = A K_t^\alpha L_t^{1-\alpha}, \quad \alpha \in (0,1),
\]

где \(A > 0\) – уровень TFP. Фирма максимизирует прибыль в каждом периоде:

\[
\max_{K_t, L_t} \; A K_t^\alpha L_t^{1-\alpha} - w_t L_t - R_t K_t.
\]

В силу совершенной конкуренции и отсутствия издержек на корректировку, условия первого порядка дают стандартные выражения для факторных цен:

\[
R_t = \alpha \frac{Y_t}{K_t}, \qquad w_t = (1-\alpha) \frac{Y_t}{L_t}.
\]

\section{Условия равновесия}

Равновесие в модели определяется как набор последовательностей \(\{C_t, L_t, K_{t+1}, Y_t, R_t, w_t\}_{t=0}^\infty\), удовлетворяющих:

\begin{itemize}
  \item задаче домохозяйства (условия оптимальности),
  \item задаче фирмы (предельные производительности факторов),
  \item бюджетному ограничению домохозяйства (которое с учётом закона движения капитала превращается в ресурсное ограничение экономики):
  \[
  K_{t+1} = Y_t - C_t + (1-\delta)K_t,
  \]
  \item начальному условию \(K_0\) (задано экзогенно)
\end{itemize}

\section{Необходимое условие экстремума}

Выпишем лагранжиан задачи домохозяйства:

\[
\mathcal{L} = \sum_{t=0}^\infty \beta^t \left[ \frac{C_t^{1-\sigma}-1}{1-\sigma} - \varphi \frac{L_t^{1+\psi}}{1+\psi} + \lambda_t \bigl( w_t L_t + R_t K_t - C_t - K_{t+1} + (1-\delta)K_t \bigr) \right],
\]

где \(\lambda_t\) – множитель Лагранжа для бюджетного ограничения в периоде \(t\).

Условия первого порядка (по \(C_t\), \(L_t\), \(K_{t+1}\)):

\begin{align}
\frac{\partial \mathcal{L}}{\partial C_t} &= 0 \quad \Rightarrow \quad \beta^t C_t^{-\sigma} - \lambda_t = 0, \tag{1}\\[4pt]
\frac{\partial \mathcal{L}}{\partial L_t} &= 0 \quad \Rightarrow \quad -\beta^t \varphi L_t^{\psi} + \lambda_t w_t = 0, \tag{2}\\[4pt]
\frac{\partial \mathcal{L}}{\partial K_{t+1}} &= 0 \quad \Rightarrow \quad -\lambda_t + \lambda_{t+1}\bigl(R_{t+1} + 1 - \delta\bigr) = 0. \tag{3}
\end{align}

Из (1) и (3) получаем уравнение Эйлера для потребления:

\[
C_t^{-\sigma} = \beta \, C_{t+1}^{-\sigma} \bigl( R_{t+1} + 1 - \delta \bigr).
\]

С учётом предельной производительности капитала \(R_{t+1} = \alpha Y_{t+1}/K_{t+1}\) уравнение принимает вид:

\[
\frac{C_{t+1}}{C_t} = \beta^{1/\sigma} \bigl( \alpha A K_{t+1}^{\alpha-1} L_{t+1}^{1-\alpha} + 1 - \delta \bigr)^{1/\sigma}.
\]

Из (1) и (2) выводится условие внутридневного выбора между потреблением и трудом (условие оптимального предложения труда):

\[
\varphi L_t^{\psi} = w_t \, C_t^{-\sigma}.
\]

Используя выражение для заработной платы \(w_t = (1-\alpha) A K_t^\alpha L_t^{-\alpha}\), получаем:

\[
\varphi L_t^{\psi} = (1-\alpha) A \left(\frac{K_t}{L_t}\right)^\alpha C_t^{-\sigma}.
\]

Таким образом, система условий первого порядка совместно с ресурсным ограничением и законом движения капитала полностью описывает динамику модели.

\section{Стационарное состояние (Steady State)}

  Стационарным состоянием (steady state) модели называется такая траектория, на которой все переменные остаются постоянными во времени. Иными словами, выполняются условия:
  \[
  K_{t+1}=K_t=K^*, \qquad 
  C_{t+1}=C_t=C^*, \qquad 
  L_{t+1}=L_t=L^*,
  \]
  и, как следствие, $Y_{t+1}=Y_t=Y^*$, $R_{t+1}=R_t=R^*$, $w_{t+1}=w_t=w^*$, $I_{t+1}=I_t=I^*$.

  Вывод аналитических выражений для стационарных значений осуществляется последовательно, начиная с уравнения Эйлера.


  Уравнение Эйлера для потребления в стационарном режиме принимает вид:
  \begin{equation}
  1 = \beta \bigl(1 + R^* - \delta\bigr).
  \end{equation}
  Отсюда немедленно получаем:
  \begin{equation}
  R^* = \frac{1}{\beta} - 1 + \delta. \tag{3.1}
  \end{equation}
  Таким образом, стационарная доходность капитала полностью определяется параметрами дисконтирования $\beta$ и выбытия $\delta$.


  Закон движения капитала в стационарном состоянии даёт:
  \[
  K^* = Y^* - C^* + (1-\delta)K^* \quad\Longrightarrow\quad Y^* = C^* + \delta K^*.
  \]

  С другой стороны, из определения доходности капитала $R^* = \alpha Y^*/K^*$ имеем $Y^* = \frac{R^*}{\alpha}K^*$. Приравнивая два выражения для $Y^*$, находим:
  \[
  \frac{R^*}{\alpha}K^* = C^* + \delta K^* \quad\Longrightarrow\quad 
  \frac{R^*}{\alpha} = \frac{C^*}{K^*} + \delta.
  \]

  Введём обозначение $v = C^*/K^*$ для отношения стационарного потребления к капиталу. Тогда:
  \begin{equation}
  v = \frac{C^*}{K^*} = \frac{R^*}{\alpha} - \delta. \tag{3.2}
  \end{equation}


  Для нахождения $L^*$ используем условие внутрипериодной оптимальности, связывающее заработную плату, потребление и труд:
  \[
  \phi C^* (L^*)^{\psi} = w^*,
  \]
  а также выражение для заработной платы через предельный продукт труда:
  \[
  w^* = (1-\alpha)\frac{Y^*}{L^*}.
  \]

  Подставляя $Y^* = \frac{R^*}{\alpha}K^*$ и $C^* = v K^*$, получаем:
  \[
  \phi v K^* (L^*)^{\psi} = (1-\alpha)\frac{R^*}{\alpha}\frac{K^*}{L^*}.
  \]

  Сокращая $K^*$ (положительный множитель) и умножая обе части на $L^*$, приходим к уравнению:
  \[
  \phi v (L^*)^{\psi+1} = \frac{1-\alpha}{\alpha}R^*.
  \]

  Отсюда выражается стационарный уровень труда:
  \begin{equation}
  L^* = \left( \frac{(1-\alpha)R^*}{\alpha \phi v} \right)^{\frac{1}{\psi+1}}. \tag{3.3}
  \end{equation}


  Теперь найдём $K^*$, используя производственную функцию Кобба–Дугласа:
  \[
  Y^* = A (K^*)^{\alpha} (L^*)^{1-\alpha}.
  \]

  С другой стороны, из $R^* = \alpha Y^*/K^*$ имеем $Y^* = \frac{R^*}{\alpha}K^*$. Приравнивая, получаем:
  \[
  \frac{R^*}{\alpha}K^* = A (K^*)^{\alpha} (L^*)^{1-\alpha}.
  \]

  Разрешая относительно $K^*$:
  \[
  (K^*)^{1-\alpha} = \frac{\alpha A}{R^*} (L^*)^{1-\alpha}
  \quad\Longrightarrow\quad
  K^* = \left( \frac{\alpha A}{R^*} \right)^{\frac{1}{1-\alpha}} L^*.
  \]

  Более удобная форма (явно выделяющая $L^*$):
  \begin{equation}
  K^* = \left( \frac{\alpha A (L^*)^{1-\alpha}}{R^*} \right)^{\frac{1}{1-\alpha}}. \tag{3.4}
  \end{equation}

  Наконец, стационарное потребление находится по определению $v$:
  \begin{equation}
  C^* = v K^*. \tag{3.5}
  \end{equation}


  Таким образом, стационарное состояние модели полностью описывается следующей системой аналитических выражений, вычисляемых в строгой последовательности:

  \[
  \begin{aligned}
  R^* &= \frac{1}{\beta} - 1 + \delta, \\[4pt]
  v &= \frac{R^*}{\alpha} - \delta, \\[4pt]
  L^* &= \left( \frac{(1-\alpha)R^*}{\alpha \phi v} \right)^{\frac{1}{\psi+1}}, \\[4pt]
  K^* &= \left( \frac{\alpha A (L^*)^{1-\alpha}}{R^*} \right)^{\frac{1}{1-\alpha}}, \\[4pt]
  C^* &= v K^*, \\[4pt]
  Y^* &= A (K^*)^{\alpha} (L^*)^{1-\alpha} = \frac{R^*}{\alpha} K^*, \\[4pt]
  I^* &= \delta K^*, \\[4pt]
  w^* &= (1-\alpha)\frac{Y^*}{L^*}.
  \end{aligned}
  \]

  Все величины выражаются исключительно через экзогенные параметры модели $\beta$, $\alpha$, $\delta$, $\phi$, $\psi$, $A$, что позволяет проводить численные расчёты и анализировать чувствительность равновесия к изменению параметров.
  